\section{VakilBuddy: A Citation-Grounded Workspace}
\label{sec:system}
VakilBuddy is a prototype litigation workspace for advocates. It is a web application organised around a \emph{matter} (a case file) whose tabs cover research, evidence, arguments, drafting, hearing preparation and case status. A matter's data stays in the user's browser; a small server serves the page and calls the model provider, and a separate service over the index answers searches and lookups. Four properties follow from the findings above.

\textbf{The index is the source of truth for authorities.} Search over the index (exact lookup by CNR, neutral or S.C.R. citation, full-text and typo-tolerant search, filters by court, year and outcome) sits in the matter's research tab, where an advocate pins the judgments the matter relies on, each with a note of what it is relied on for. A citation check finds the citations and case names in a draft or in pasted text with patterns, looks each one up in the index, and reports whether the case exists, whether the number belongs to another case, whether the number is wrong, whether nothing is found, or whether it is a High Court case that cannot be checked (\S\ref{sec:case}). This check is simpler than the verifier of \S\ref{sec:verifier}: it expands state abbreviations and tolerates a one-letter spelling slip, but it searches the index by whole words, so a title that the archive spells differently is not found by name. Like the verifier it confirms existence, not that a judgment supports the proposition.

\textbf{Model text is bounded by sources.} Model requests go to OpenRouter \cite{openrouter} or any OpenAI-compatible endpoint; the key is kept on the server and never returned to the browser. A draft written by the model may cite only the authorities pinned to the matter and list only the matter's own documents as annexures, leaves a visible blank where a fact is missing, and has its citations checked in the same response. Every model answer carries the model's name, and where an action names cases they are looked up in the index and marked. For the actions that read the advocate's own documents, an answer is kept only if the documents support it: a position attributed to a witness is kept only if its content words, and every date and number in it, appear in that witness's statement; a next hearing date is kept only if the order writes it; and a chronology's source, and the document a witness is to be confronted with, must be documents that were supplied, with a dated event kept only if the documents write its date. Items that fail are dropped and counted. A failed model call is shown as an error and never replaced by a canned answer, and the percentages a model may offer for risk are not used.

\textbf{Evidence is read where it lies.} Text, Word and PDF files with a text layer are read directly, with pdf.js \cite{pdfjs}; scanned pages and photographs are read by Tesseract \cite{smith2007} in English, Hindi and Marathi inside the browser, so the files are never uploaded and only the text is kept. The text is searched across a matter's documents. Dates and rupee amounts, including Devanagari digits, are found by pattern together with the nearest event word, and two documents that give different dates or amounts beside the same word are listed as possible inconsistencies: places to look, never findings.

\textbf{Dates and drafts follow rules the advocate can inspect.} A limitation calculator holds \CasePeriods{} periods, each naming its source (the Limitation Act, 1963 \cite{limitationact}, the Supreme Court Rules, 2013 \cite{scrules}, the Code of Civil Procedure and others). It leaves out the day of the event (s.~12(1)) and, where the law allows it, the time taken to obtain a certified copy (s.~12(2)), runs a period that ends on a closed day to the day the court reopens (s.~4), and shows how each date was counted. Drafting offers \CaseFormats{} built-in formats, keeps the last twelve versions of a draft with a word-level redline between any two, takes reviewers' comments, and exports a Word file only after the advocate approves. The workspace files and sends nothing; a WhatsApp reminder opens with its text written and the advocate presses Send.

\textbf{Status.} The workspace is covered by \CaseTestsWorkspace{} automated tests and its index service by \CaseTestsApi{}. Model calls have been tested against a stub provider and a local open-weight model; hosted models, including those reached through OpenRouter, are untested. Apart from the matter of \S\ref{sec:case}, the matters in the workspace are sample data, and there is no user study.
